% TOP_PROBLEM: {"id": 30005163, "title": "Fourier Restriction Conjecture in Higher Dimensions", "category_id": 8, "category": {"id": 8, "name": "analysis", "display_name": "Analysis", "description": "Limits, continuity, calculus, and function theory.", "slug": "analysis", "order_index": 8, "created_at": "2026-07-31T15:26:25.671Z"}, "status": "open"}
% FIRST_POSTED: 2026-10-06
% ATTEMPT_STATUS: unresolved
% OPENAI_CLAIM: {"schema_version":1,"source_repo":"https://github.com/openai/math","source_commit":"adc7f1241b42e322a6451854ab7e4b4c146bf78a","release_date":"2026-10-06","match":"partial","resolution":"partial","independently_reviewed":false,"second_pass":"n/a","families":[{"family":"077","title":"Fourier restriction for positively curved surfaces","subject":"Real and complex analysis","manuscripts":[{"title":"Diagonal Fourier extension for positively curved surfaces in three dimensions","pdf_path":"preprints/Diagonal-Fourier-extension-for-positively-curved-surfaces-in-three-dimensions-September-24-2026/Diagonal-Fourier-extension-for-positively-curved-surfaces-in-three-dimensions-September-24-2026.pdf","readme_path":"preprints/Diagonal-Fourier-extension-for-positively-curved-surfaces-in-three-dimensions-September-24-2026/README.md","date":"2026-09-24","theorem_ref":"thm:main (sections/00-introduction.tex:18)","statement_sources":[{"path":"preprints/Diagonal-Fourier-extension-for-positively-curved-surfaces-in-three-dimensions-September-24-2026/build/sections/00-introduction.tex","line":18,"label":"thm:main"}]},{"title":"Elliptic capacity propagation and Fourier restriction to the sphere","pdf_path":"preprints/Elliptic-capacity-propagation-and-Fourier-restriction-to-the-sphere-September-24-2026/Elliptic-capacity-propagation-and-Fourier-restriction-to-the-sphere-September-24-2026.pdf","readme_path":"preprints/Elliptic-capacity-propagation-and-Fourier-restriction-to-the-sphere-September-24-2026/README.md","date":"2026-09-24","theorem_ref":"thm:main (sections/00-introduction.tex:21)","statement_sources":[{"path":"preprints/Elliptic-capacity-propagation-and-Fourier-restriction-to-the-sphere-September-24-2026/build/sections/00-introduction.tex","line":21,"label":"thm:main"}]}],"lean":null,"reasoning_trace":null}]}
\documentclass[11pt]{article}
\usepackage[margin=25mm]{geometry}
\usepackage{fontspec}
\setmainfont{Latin Modern Roman}
\usepackage{amsmath,amssymb,mathtools}
\usepackage{unicode-math}
\setmathfont{Latin Modern Math}
\usepackage{enumitem}
\usepackage{xurl,hyperref,bookmark}
\hypersetup{colorlinks=true,urlcolor=blue}
\setcounter{secnumdepth}{0}
\setlength{\parindent}{0pt}
\setlength{\parskip}{5pt}
\setlength{\emergencystretch}{3em}
\begin{document}
\section*{Fourier Restriction Conjecture in Higher Dimensions}

\subsection{OpenAI's claim}
\textit{Extension from curved surfaces}: For a compact smooth surface in $\mathbb R^3$ with positive curvature, the paper claims diagonal Fourier extension boundedness for every exponent $p>3$; it includes smooth surface boundaries and draws a two-dimensional Schrödinger smoothing consequence from the paraboloid case.

\textit{Sphere extension and packet capacity}: For bounded data on the two-sphere, the announced extension bound has every output exponent larger than three; the threshold three is also asserted to be the sharp endpoint of this open range.

\subsection{Mathematical statement of the claimed result}
\paragraph{Mathematical claim 1.} Let $\Sigma\subset\mathbb R^3$ be a compact smooth surface with positive curvature, allowing smooth boundary, and use surface area $\sigma$. For $E_\Sigma f(x)=\int_\Sigma e^{2\pi ix\cdot\xi}f(\xi)\,d\sigma(\xi)$, the claim is \[\|E_\Sigma f\|_{L^p(\mathbb R^3)}\le C_{\Sigma,p}\|f\|_{L^p(\Sigma,\sigma)}\quad(3<p<\infty),\] for all complex-valued $f\in L^p(\Sigma)$. It concerns this diagonal range in three dimensions, with a surface-dependent constant.

\paragraph{Mathematical claim 2.} For surface area $\sigma$ on $S^2$ and $Eg(x)=\int_{S^2}e^{2\pi ix\cdot\omega}g(\omega)\,d\sigma(\omega)$, every bounded measurable complex $g$ satisfies \[\|Eg\|_{L^p(\mathbb R^3)}\le C_p\|g\|_{L^\infty(S^2,\sigma)}\quad(p>3).\] The constant depends only on $p$; the endpoint $p=3$ and other ambient dimensions are outside this theorem.

\subsection{Outline of the argument}
\textit{Extension from curved surfaces}: Graph patches are extended to phases compatible with an elliptic-capacity propagation theorem. That theorem is combined with weighted Kakeya control to sum packet coefficients by their capacity. A joint estimate for widely separated balls removes the local spatial loss, and chart assembly transfers the estimate to the original surface.

\textit{Sphere extension and packet capacity}: Joint tests of position and velocity use ellipse weights that penalize eccentricity. A geometric propagation estimate controls how this capacity evolves along lines, using overlap savings from differently directed tests. Exact packet analysis and synthesis retain the oscillatory information needed to convert the positive capacity estimate into Fourier extension bounds.

\subsection{Scope relative to this problem}
\textit{Extension from curved surfaces}: The extension estimate establishes the stated p>3 branch in ambient dimension three for spheres or positively curved surfaces. The question asks the sharp two-exponent restriction range in all dimensions or for general submanifolds, which is broader in dimensions, surfaces, and/or exponents.

\textit{Sphere extension and packet capacity}: The extension estimate establishes the stated p>3 branch in ambient dimension three for spheres or positively curved surfaces. The question asks the sharp two-exponent restriction range in all dimensions or for general submanifolds, which is broader in dimensions, surfaces, and/or exponents.

This page records partial support. Complete resolution of the catalogue question is not established by this record.

\subsection{Formal verification}
\textit{Extension from curved surfaces}: The pinned catalogue lists no Lean formalization for this result.

\textit{Sphere extension and packet capacity}: The pinned catalogue lists no Lean formalization for this result.

This site has not run the Lean development. These links describe OpenAI’s selected formal statements and should be read with their stated scope.

\subsection{Status caveat}
OpenAI reports that an internal model produced the manuscripts. Their arguments have not been independently verified as proofs by this site. Any portion outside the documented formal scope may contain mistakes, and the catalogue label reports the announced mathematical claim.

\subsection{Sources}
\begin{itemize}
\item OpenAI, \emph{Diagonal Fourier extension for positively curved surfaces in three dimensions}, 2026-09-24; citation key \texttt{OAI:Diagonal-Fourier-extension-for-positively-curved-surfaces-in-three-dimensions-September-24-2026}. \href{https://github.com/openai/math/blob/adc7f1241b42e322a6451854ab7e4b4c146bf78a/preprints/Diagonal-Fourier-extension-for-positively-curved-surfaces-in-three-dimensions-September-24-2026/Diagonal-Fourier-extension-for-positively-curved-surfaces-in-three-dimensions-September-24-2026.pdf}{PDF (GitHub)}; \href{https://github.com/openai/math/raw/adc7f1241b42e322a6451854ab7e4b4c146bf78a/preprints/Diagonal-Fourier-extension-for-positively-curved-surfaces-in-three-dimensions-September-24-2026/Diagonal-Fourier-extension-for-positively-curved-surfaces-in-three-dimensions-September-24-2026.pdf}{PDF (direct)}; \href{https://github.com/openai/math/blob/adc7f1241b42e322a6451854ab7e4b4c146bf78a/preprints/Diagonal-Fourier-extension-for-positively-curved-surfaces-in-three-dimensions-September-24-2026/README.md}{Manuscript page and citation}; \href{https://github.com/openai/math/blob/main/preprints/Diagonal-Fourier-extension-for-positively-curved-surfaces-in-three-dimensions-September-24-2026/Diagonal-Fourier-extension-for-positively-curved-surfaces-in-three-dimensions-September-24-2026.pdf}{Latest version}; Statement source: \href{https://github.com/openai/math/blob/adc7f1241b42e322a6451854ab7e4b4c146bf78a/preprints/Diagonal-Fourier-extension-for-positively-curved-surfaces-in-three-dimensions-September-24-2026/build/sections/00-introduction.tex\#L18}{thm:main}.
\item OpenAI, \emph{Elliptic capacity propagation and Fourier restriction to the sphere}, 2026-09-24; citation key \texttt{OAI:Elliptic-capacity-propagation-and-Fourier-restriction-to-the-sphere-September-24-2026}. \href{https://github.com/openai/math/blob/adc7f1241b42e322a6451854ab7e4b4c146bf78a/preprints/Elliptic-capacity-propagation-and-Fourier-restriction-to-the-sphere-September-24-2026/Elliptic-capacity-propagation-and-Fourier-restriction-to-the-sphere-September-24-2026.pdf}{PDF (GitHub)}; \href{https://github.com/openai/math/raw/adc7f1241b42e322a6451854ab7e4b4c146bf78a/preprints/Elliptic-capacity-propagation-and-Fourier-restriction-to-the-sphere-September-24-2026/Elliptic-capacity-propagation-and-Fourier-restriction-to-the-sphere-September-24-2026.pdf}{PDF (direct)}; \href{https://github.com/openai/math/blob/adc7f1241b42e322a6451854ab7e4b4c146bf78a/preprints/Elliptic-capacity-propagation-and-Fourier-restriction-to-the-sphere-September-24-2026/README.md}{Manuscript page and citation}; \href{https://github.com/openai/math/blob/main/preprints/Elliptic-capacity-propagation-and-Fourier-restriction-to-the-sphere-September-24-2026/Elliptic-capacity-propagation-and-Fourier-restriction-to-the-sphere-September-24-2026.pdf}{Latest version}; Statement source: \href{https://github.com/openai/math/blob/adc7f1241b42e322a6451854ab7e4b4c146bf78a/preprints/Elliptic-capacity-propagation-and-Fourier-restriction-to-the-sphere-September-24-2026/build/sections/00-introduction.tex\#L21}{thm:main}.
\item \href{https://github.com/openai/math/blob/adc7f1241b42e322a6451854ab7e4b4c146bf78a/CONTENTS.md}{OpenAI catalogue entry}.
\end{itemize}
\end{document}
